\Needspace{6\baselineskip}
\begin{center}\hypertarget{native-261-article188}{}\pdfbookmark[2]{188.}{article-188}188.\end{center}
The preceding method is chiefly suited to cases where the highest precision
is desired. Very often, however, one may readily give up a little precision
if this considerably shortens the calculation, especially when the observations
do not yet cover a long interval and one is therefore not yet aiming at what
might be called a definitive determination of the orbit. In such cases the
following method can be used to considerable advantage.

From the full set of observations, select two complete positions $L$ and $L'$
and calculate the distances of the celestial body from the Earth at the
corresponding times from the approximate elements. Then form three hypotheses
for these distances. Retain the calculated values in the first, change the
first distance in the second hypothesis, and change the second distance in
the third. Each change may be chosen at will in proportion to the uncertainty
presumed to remain in the corresponding distance. Under the three hypotheses
shown in the following table,
\begin{center}
\begingroup
\setlength{\tabcolsep}{4pt}
\renewcommand{\arraystretch}{1.3}
\begin{tabular}{|l|c|c|c|}
\hline
 & Hyp. I & Hyp. II & Hyp. III\\
\hline
Distance\footnotemark\ corresponding to the first position & $D$ & $D+\delta$ & $D$\\
Distance corresponding to the second position & $D'$ & $D'$ & $D'+\delta'$\\
\hline
\end{tabular}
\endgroup
\end{center}
\footnotetext{It is even more convenient to use the logarithms of the curtate distances in place of the distances themselves.}
\noindent\SourceLeaf{262}calculate three sets of elements from the two
positions $L,L'$ by the methods explained in Book One. Then calculate from
each set the geocentric positions of the celestial body corresponding to the
times of all the remaining observations. Denote them as follows, with each
longitude and latitude, or right ascension and declination, denoted separately.
\begin{center}
\begin{tabular}{@{}l@{\quad}l@{\quad}l@{\quad}l@{\quad}l@{}}
\makebox[12.8em][l]{in the first set\hfill .\,.\,.\,.\,.\,.} & $M$, & $M'$, & $M''$ & etc.\\[.3em]
\makebox[12.8em][l]{in the second set\hfill .\,.\,.\,.\,.} & $M+a$, & $M'+a'$, & $M''+a''$ & etc.\\[.3em]
\makebox[12.8em][l]{in the third set\hfill .\,.\,.\,.\,.\,.\,.} & $M+\beta$, & $M'+\beta'$, & $M''+\beta''$ & etc.
\end{tabular}
\end{center}
\noindent Further, let the corresponding observed positions be
\begin{center}
\begin{tabular}{@{}l@{\quad}l@{\quad}l@{\quad}l@{\quad}l@{}}
\makebox[12.8em][l]{observed positions\hfill .\,.\,.\,.\,.\,.\,.\,.\,.\,.\,.} & $N$, & $N'$, & $N''$ & etc.
\end{tabular}
\end{center}

Now, insofar as small changes in the distances $D,D'$ produce proportional
changes in the individual elements and in the geocentric positions calculated
from them, we may assume that the geocentric positions calculated from a
fourth set of elements based on the distances $D+x\delta,D'+y\delta'$ from
the Earth will be $M+ax+\beta y,M'+a'x+\beta'y,M''+a''x+\beta''y$ etc.,
respectively. Following the preceding investigations, determine $x,y$ so
that these quantities agree as closely as possible with the corresponding
$N,N',N''$ etc., taking into account the relative precision of the observations.
The corrected set of elements itself can be derived in the same way from
$L,L'$ and the distances $D+x\delta,D'+y\delta'$, or by simple interpolation
from the first three sets of elements according to the known rules.
